% ============================================================
% APPENDIX B -- General Dynamical Extension of the H-potential
% Rewritten to mirror Appendix A: Purpose -> Definitions ->
% Propositions -> Proofs -> Remarks -> Summary.
% ============================================================

\label{app:dynamical}

\begin{ataglance}[How to read this appendix]
Appendix~\ref{app:mathematical_foundations} asked what a person's health
capability is at one moment. This appendix asks what it does next, and it
answers in three layers of increasing realism.

The first layer is a single ordinary differential equation for one patient with
no randomness (\S\ref{sec:B_ode}). It says the disease marker rises with
treatment and repair and falls with inhibition and clearance, and it settles at
a level where those balance. The second layer adds noise, because two patients
on identical regimens do not follow identical curves (\S\ref{sec:B_sde}). The
third layer stops following individuals and follows the whole cohort as a
density that flows and spreads (\S\ref{sec:B_fp}). The last section reconnects
the machinery to the conceptual score by letting it move in time
(\S\ref{sec:B_timevarying}).

The clinical reader may recognise the whole structure from pharmacokinetics.
Layer one is a one-compartment model; layer two is that model with
between-dose variability; layer three is the population-PK view in which the
object of interest is the distribution of trough levels across a clinic rather
than any one patient's curve. Profile~B of Table~\ref{tab:profiles} is the
reference case throughout, and Appendix~\ref{app:hemo_dynamics} makes the haemophilia version explicit.

None of this is used in the empirical validation. It is a proposed direction.
\end{ataglance}

\section{Purpose and Scope}
\label{sec:B_purpose}
This appendix records a disease-independent dynamical extension of the
$H_{\text{potential}}$ framework of Chapter~\ref{ch:theory} and
Appendix~\ref{app:mathematical_foundations}. Where Appendix~\ref{app:mathematical_foundations}
treats the static conceptual score and its measurement, this appendix asks how
the biomedical state that feeds that score evolves in time, and how individual
trajectories aggregate to a population distribution. The material is a proposed
direction for future work: the specifications are theoretical and the figures
are illustrative rather than empirical. Nothing here is used in the measurement
or prognostic validation of Chapters~\ref{ch:methods} to \ref{ch:results2}.

Saying what the appendix is not part of leaves open what it is attached to, and
the mapping between the two matters for reading the notation.

\section{Relationship to the Validated Framework}
\label{sec:B_relationship}
The definitions of the conceptual variables $L,Q,S,B$ and the latent construct
$\zeta$ are given in Appendix~\ref{app:mathematical_foundations}. This appendix
adds a state variable $X(t)$ that stands behind the biomedical term $L$: in the
time-varying form of Section~\ref{sec:B_timevarying}, $L$ becomes a function of
$X(t)$, so the static score of Appendix~\ref{app:mathematical_foundations} is the
$t$-independent special case of the construction here. The disease-specific
realisation of this appendix is Appendix~\ref{app:hemo_dynamics}, and its
numerical implementation is Appendix~\ref{app:computation}.

Each proposition below is stated as a claim a computer can check, and each is
checked. The line beneath a proposition names the test that asserts it; the
tests live in \texttt{verification/tests/test\_dynamics.py} in the thesis
archive and Appendix~\ref{app:verification} reports the outcome of running
them.

With the correspondence fixed, the objects can be defined precisely enough to
prove things about.

\section{Definitions}
\label{sec:B_definitions}

\begin{definition}[Deterministic state]
\label{def:B_state}
Let $X(t)\in\mathbb{R}_+$ denote the latent disease state of an individual,
measured by an observable biomarker appropriate to the disease (for example
factor~VIII concentration, viral load, tumour burden, or HbA1c). The remaining
mechanism variables are the inhibitory factor $Y\ge0$, the homeostatic
(restorative) factor $Z\ge0$, the therapeutic input $T\ge0$, and the natural
clearance rate $k_{\text{elim}}>0$.
\end{definition}

\begin{definition}[Universal deterministic model]
\label{def:B_ode}
The disease marker evolves by the ordinary differential equation \citep{coddington_levinson_1955}
\begin{equation}
\frac{dX}{dt}=\mathbf{f}(T)-\mathbf{g}(Y,X)+\mathbf{h}(Z)-k_{\text{elim}}X,
\label{eq:B_ode}
\end{equation}
with treatment input $\mathbf{f}(T)=V_{\max}T/(K_m+T)$ (Michaelis--Menten),
cooperative stabilisation $\mathbf{h}(Z)=\gamma Z^n/(K_d^n+Z^n)$ (Hill), with
maximal stabilisation rate $\gamma>0$, and bilinear
inhibition $\mathbf{g}(Y,X)=k_{\text{inh}}YX$.
\end{definition}

\begin{definition}[Patient state and stochastic model]
\label{def:B_sde}
The patient state is the stochastic process $X(t)$ governed by \citep{oksendal_2003}
\begin{equation}
dX(t)=\big[\mathbf{f}-\mathbf{g}+\mathbf{h}-k_{\text{elim}}X(t)\big]\,dt+\sigma(X,t)\,dW_t,
\label{eq:B_sde}
\end{equation}
where $W_t$ is a standard Wiener process \citep{oksendal_2003} and $\sigma(X,t)\ge0$ scales the
perturbation. Patient-specific treatment sensitivity $S_T=f(\{G_i\},\{I_j\},\text{env})$
modulates the drift, personalising \eqref{eq:B_sde}.
\end{definition}

\begin{definition}[Population state]
\label{def:B_population}
The population state is the probability density $\rho(x,t)$ of patient states
$X(t)$ across a cohort, with drift $\mu(x,t)$ inherited from the SDE and
diffusion $D(x,t)=\sigma(x,t)\sigma(x,t)^\top$.

Three objects have now been defined: a patient, a noisy patient and a
population. They are taken in that order, because each is built from the one
before and because the first is the only one for which existence can be settled
by hand.
\end{definition}

\section{The Deterministic Backbone}
\label{sec:B_ode}

\begin{proposition}[Existence and uniqueness]
\label{prop:B_existence}
If the parameters are bounded and $\mathbf{f},\mathbf{g},\mathbf{h}$ are
Lipschitz \citep{coddington_levinson_1955} on the admissible range, then the right-hand side of \eqref{eq:B_ode}
is globally Lipschitz in $X$ with constant $L=k_{\text{inh}}Y+k_{\text{elim}}$, and
\eqref{eq:B_ode} admits a unique solution for each initial condition. Moreover,
since the $X$-derivative of the right-hand side equals $-L<0$, two solutions
whose initial conditions differ by $\delta_0$ satisfy
$|X_1(t)-X_2(t)|=\delta_0e^{-Lt}$: trajectories contract, and contract faster
the larger $L$ is.
\end{proposition}
\begin{proof}
Only $-\mathbf{g}(Y,X)-k_{\text{elim}}X=-(k_{\text{inh}}Y+k_{\text{elim}})X$
depends on $X$ (the saturable $\mathbf{f}$ and $\mathbf{h}$ do not), so the
$X$-derivative of the right-hand side is the constant
$-(k_{\text{inh}}Y+k_{\text{elim}})$; hence the map is globally Lipschitz with
that constant. Picard--Lindel\"of then gives local existence and uniqueness, and
the linear growth bound extends the solution to all $t\ge0$. Because the
right-hand side depends on $X$ only through the term $-LX$, the difference
$\delta(t)$ of two solutions satisfies $\dot\delta=-L\delta$, so
$\delta(t)=\delta_0e^{-Lt}$; Gr\"onwall's inequality \citep{gronwall_1919}
gives the same bound when $\mathbf{f}$ and $\mathbf{h}$ are only Lipschitz.
\end{proof}

\begin{reading}[Why an existence proof is not a formality]
This proposition says that the model has exactly one answer and that the answer
does not jump. Both halves matter clinically.

Uniqueness means that if two people run the simulation with the same starting
factor level and the same regimen, they get the same trajectory. A model
without it can produce two different predictions from one set of inputs, which
would make any dosing recommendation drawn from it arbitrary.

The Gr\"onwall bound is the stability half. It says that if the true
starting factor activity differs from what the assay reported, the gap between
the two predicted trajectories is bounded rather than growing without limit.
Since every clinical input is measured with error, a model lacking this
property would amplify assay noise into confident nonsense.

Here the bound is generous, because the $X$-derivative of the right-hand side
is $-(k_{\text{inh}}Y+k_{\text{elim}})$, which is negative. Two solutions
starting $\delta_0$ apart therefore close to $\delta_0e^{-Lt}$: they converge,
and they converge faster the larger $L$ is. An error in the assumed starting
factor level is forgotten, and a higher inhibitor titre makes it forgotten
sooner, because neutralisation and clearance together pull every trajectory
towards the same steady state. What a high titre does not do is make the
steady state itself easy to hit, which is the subject of the next
proposition.
\end{reading}

\begin{workedexample}[apb-gronwall]{Two trajectories from two assay readings}
Let the true starting factor activity be $2$ and the assay reading $3$, so
$\delta_0=1$. Take the illustrative parameter set of
\texttt{verification/tests/test\_dynamics.py}: $V_{\max}=100$, $K_m=10$, $T=20$,
$k_{\text{inh}}=0.05$, $Y=2$, $\gamma=30$, $K_d=50$, $n=2$, $Z=80$,
$k_{\text{elim}}=0.8$. The Lipschitz constant is
$L=k_{\text{inh}}Y+k_{\text{elim}}=0.05\times2+0.8=0.9$.

Integrating both starting values and taking the difference gives $0.407$ at
$t=1$, $0.011$ at $t=5$ and $0.00012$ at $t=10$, matching $e^{-0.9t}$ to three
figures at every point. Raise the titre to $Y=20$ and $L$ becomes $1.8$: the
gap at $t=5$ falls to $0.00012$, ninety times smaller than at $Y=2$.

\begin{wecode}[Python]
import math
from hpv.dynamics import integrate_ode, lipschitz_constant

P = dict(Vmax=100.0, Km=10.0, T=20.0, k_inh=0.05, Y=2.0,
         gamma=30.0, Kd=50.0, n=2.0, Z=80.0, k_elim=0.8)
L = lipschitz_constant(k_inh=P["k_inh"], Y=P["Y"], k_elim=P["k_elim"])
print("L =", L)
for t in (1.0, 2.0, 5.0, 10.0):
    gap = abs(integrate_ode(2.0, t, P) - integrate_ode(3.0, t, P))
    print(f"t {t:>5.1f}   gap {gap:.6f}   exp(-Lt) {math.exp(-L * t):.6f}")

hi = dict(P, Y=20.0)
print("gap at t=5 with Y=20:",
      round(abs(integrate_ode(2.0, 5.0, hi) - integrate_ode(3.0, 5.0, hi)), 6))
\end{wecode}

\weanswer{$L=0.9$; the gap falls from $0.407$ at $t=1$ to $1.23\times10^{-4}$ at
$t=10$, tracking $e^{-Lt}$ throughout. At $Y=20$ the same gap is reached by
$t=5$. Assay error is forgotten, and forgotten sooner at a higher titre.
Asserted by
\texttt{test\_worked\_example\_gronwall\_two\_assay\_readings} and
\texttt{test\_proposition\_trajectories\_converge\_not\_diverge}.}
\end{workedexample}
\checkedby{tests/test_dynamics.py::test_proposition_existence_lipschitz_constant}

\begin{proposition}[Deterministic equilibrium]
\label{prop:B_equilibrium}
For constant inputs the ODE has the unique stable fixed point
$X^\ast=(\mathbf{f}+\mathbf{h})/(k_{\text{inh}}Y+k_{\text{elim}})$.
\end{proposition}
\begin{proof}
Setting $dX/dt=0$ gives
$\mathbf{f}+\mathbf{h}=(k_{\text{inh}}Y+k_{\text{elim}})X^\ast$; the Jacobian
$-(k_{\text{inh}}Y+k_{\text{elim}})<0$ makes it stable.
\end{proof}

\begin{reading}[The steady state as a trough level]
$X^\ast=(\mathbf f+\mathbf h)/(k_{\text{inh}}Y+k_{\text{elim}})$ is input over
output: everything putting factor into the circulation, divided by everything
taking it out. It is the same arithmetic as a steady-state drug concentration,
dose rate over clearance, and it can be read the same way.

For Profile~B the numerator is the infused factor and whatever endogenous
stabilisation von~Willebrand factor provides. The denominator has two parts, and only one is
ordinary. Clearance $k_{\text{elim}}$ is the usual half-life term. The other,
$k_{\text{inh}}Y$, is the inhibitor, and it enters the denominator, so anything
that doubles the denominator halves the steady state for a given dose.

That is why the arithmetic of inhibitors is so unforgiving. Raising the dose
raises the numerator linearly, so it can compensate, but only in proportion,
and the ceiling is set by what can be infused and afforded. The equation makes
the case for immune tolerance induction quantitative: reducing $Y$ acts on the
denominator, and is worth more per unit than pushing the numerator.
\end{reading}

\begin{workedexample}[apb-trough]{The steady state read as a trough level}
Use the same illustrative parameters as Worked Example~\ref{ex:apb-gronwall}.
The infusion term is $\mathbf f=V_{\max}T/(K_m+T)=100\times20/30=66.67$ and the
stabilisation term is $\mathbf h=\gamma Z^n/(K_d^n+Z^n)=30\times
6400/8900=21.57$, so the numerator is $88.24$. The denominator is
$k_{\text{inh}}Y+k_{\text{elim}}=0.9$. The trough is therefore
$X^\ast=88.24/0.9=98.04$.

Now raise the titre from $Y=2$ to $Y=20$. The denominator becomes
$0.05\times20+0.8=1.8$, exactly twice as large, and the trough falls to
$49.02$, exactly half. It takes a tenfold rise in titre to halve the trough
here only because clearance still dominates the denominator at $Y=2$; in a
patient whose clearance term were smaller, the same halving would follow a much
smaller rise.

\begin{wecode}[Python]
from hpv.dynamics import equilibrium, michaelis_menten, hill, integrate_ode

P = dict(Vmax=100.0, Km=10.0, T=20.0, k_inh=0.05, Y=2.0,
         gamma=30.0, Kd=50.0, n=2.0, Z=80.0, k_elim=0.8)
f = michaelis_menten(P["T"], P["Vmax"], P["Km"])
h = P["gamma"] * hill(P["Z"], P["Kd"], P["n"])
den = P["k_inh"] * P["Y"] + P["k_elim"]
print(f"numerator {f + h:.3f} = {f:.3f} + {h:.3f}   denominator {den:.3f}")
print(f"trough {equilibrium(**P):.3f}   solver from x0=0: "
      f"{integrate_ode(0.0, 60.0, P):.3f}")
print("trough at Y=20:", round(equilibrium(**dict(P, Y=20.0)), 3))
\end{wecode}

\weanswer{$X^\ast=98.044$, which the integrator reaches from any starting
value; at $Y=20$ it is $49.022$, exactly half. Asserted by
\texttt{test\_worked\_example\_trough\_level\_for\_profile\_b} and
\texttt{test\_proposition\_equilibrium\_matches\_the\_solver}.}
\end{workedexample}

Stability also tells us what the model cannot do. It converges to one level and
stays there, so it cannot generate a spontaneous bleed in a patient whose
average factor level is adequate. Real patients have those. The next layer adds
the randomness that produces them.
\checkedby{tests/test_dynamics.py::test_proposition_equilibrium_matches_the_solver}

\section{Stochastic Extension}
\label{sec:B_sde}

\begin{proposition}[Reduction to the ODE]
\label{prop:B_reduction}
As $\sigma\to0$ the solution of the SDE \eqref{eq:B_sde} converges to the
solution of the ODE \eqref{eq:B_ode} with the same initial condition.
\end{proposition}
\begin{proof}
With $\sigma\equiv0$ the stochastic integral vanishes and \eqref{eq:B_sde}
reduces to \eqref{eq:B_ode} pathwise. For $\sigma>0$, take expectations of
\eqref{eq:B_sde} and use $\mathbb E[dW_t]=0$: the drift of \eqref{eq:B_ode} is
affine in $X$ on the admissible range, so $\mathbb E[X(t)]$ satisfies
\eqref{eq:B_ode} exactly, as Proposition~\ref{prop:B_consistency} states. The
deviation of any individual path from that mean is $O(\sigma)$ in the
mean-square sense by the standard It\^o isometry bound, so the law of $X(t)$
concentrates on the deterministic solution as $\sigma\to0$. The numerical
check reported in Appendix~\ref{app:verification} confirms the rate rather than
establishing it.
\end{proof}
\checkedby{tests/test_dynamics.py::test_proposition_sde_reduces_to_ode_as_noise_vanishes}

\begin{proposition}[Population consistency]
\label{prop:B_consistency}
If the drift is affine in $X$ and the noise has zero mean, the expectation
$\mathbb{E}[X(t)]$ satisfies the deterministic ODE. For the linear model
$dX=(a-bX)\,dt+\sigma\,dW$, $\ \mathbb{E}[X(t)]=a/b+(x_0-a/b)e^{-bt}$.
\end{proposition}
\begin{proof}
Taking expectations of \eqref{eq:B_sde} and using $\mathbb{E}[dW_t]=0$ gives
$d\mathbb{E}[X]/dt=\mathbb{E}[\text{drift}]$, which for affine drift equals the
drift evaluated at $\mathbb{E}[X]$; solving $m'=a-bm$ yields the stated mean.
\end{proof}

\begin{reading}[The average patient is not a patient]
These two propositions are a pair, and the second contains a trap worth
labelling.

The first says the stochastic model contains the deterministic one: turn the
noise down and the two coincide. That is the sanity check any extension should
pass.

The second says that when the drift is affine, the \emph{mean} of the
stochastic paths follows the deterministic curve exactly. It is tempting to
conclude that the deterministic model is therefore good enough. It is not,
because clinical outcomes are rarely affine functions of the state. What
matters for Profile~B is the proportion of hours spent below the protective
threshold, not mean factor activity, and that proportion is a highly
non-linear function of the path.

Two patients can share a mean trough of $3$~IU/dL, one holding steady and one
oscillating between $0.5$ and $5.5$. The deterministic model calls them
identical. The second bleeds. This is the same reason that HbA1c and glycaemic
variability are both measured: an average conceals the excursions that cause
the damage. Jensen's inequality is the formal name for the gap, and the whole
point of the stochastic layer is to keep it visible.
\end{reading}
\checkedby{tests/test_dynamics.py::test_proposition_population_consistency_affine_drift}

\begin{remark}
The Euler--Maruyama and Milstein schemes integrate \eqref{eq:B_sde}; the Milstein
correction $\tfrac12\sigma\sigma'(dW^2-dt)$ raises the strong order to one. The
Euler--Maruyama scheme is implemented in \texttt{hpv/dynamics.py}; the Milstein
correction is not, because none of the results below turn on the difference.
\end{remark}

Simulating enough individual paths to describe a clinic is expensive and
indirect. The density of the whole cohort obeys an equation of its own, and
solving that equation once is cheaper than following everybody.

\section{Population-Level Extension: the Fokker--Planck Equation}
\label{sec:B_fp}

\begin{definition}[Probability density, drift, diffusion]
\label{def:B_fp_terms}
Let $\rho(x,t)$ be the density of patient states, normalised so
$\int\rho(x,t)\,dx=1$; let $\mu(x,t)$ be the drift of \eqref{eq:B_sde} and
$D(x,t)=\sigma\sigma^\top$ the diffusion.
\end{definition}

The density evolves by the Fokker--Planck (Kolmogorov forward) equation \citep{risken_1996}
\begin{equation}
\frac{\partial\rho}{\partial t}
=-\frac{\partial}{\partial x}\big[\mu(x,t)\rho\big]
+\frac{1}{2}\frac{\partial^2}{\partial x^2}\big[D(x,t)\rho\big]
=-\frac{\partial J}{\partial x},\qquad
J=\mu\rho-\tfrac12\frac{\partial}{\partial x}(D\rho).
\label{eq:B_fp}
\end{equation}

\begin{proposition}[Mass conservation]
\label{prop:B_mass}
Under reflecting (zero-flux) boundaries $J=0$ at the ends of the domain, total
probability is conserved: $\tfrac{d}{dt}\int\rho\,dx=0$.
\end{proposition}
\begin{proof}
Integrating \eqref{eq:B_fp} over the domain $[x_L,x_R]$,
$\tfrac{d}{dt}\int_{x_L}^{x_R}\rho\,dx=-[J]_{x_L}^{x_R}=0$ when the boundary
fluxes vanish. The conservative finite-volume discretisation inherits this: the
assembled operator has zero column sums (numerically $\sim10^{-14}$), so the
backward-Euler step preserves total mass exactly.
\end{proof}

\begin{reading}[What the Fokker--Planck equation is actually saying]
Stop tracking patients and start tracking the crowd. Let $\rho(x,t)$ be the
proportion of a clinic population whose factor activity is near $x$ at time
$t$: a histogram that changes shape as the day goes on.

The equation says the histogram changes for exactly two reasons, and the two
terms name them. The first term, the one with the single derivative, is
\emph{drift}: everyone is being pushed the same way, so the whole histogram
slides. Start a clinic on prophylaxis and the bulk of the distribution moves
right. The second term, with the second derivative, is \emph{diffusion}:
individual variation spreads the histogram out, flattening the peak and
thickening the tails, without moving its centre.

The clinically important quantity is usually the left tail, the fraction of the
clinic below the protective threshold at any moment. A policy can shift that
fraction in two quite different ways, and the equation separates them. Raising
everyone's dose is drift. Reducing variability, through better adherence
support, more reliable supply, or a longer-acting product, is diffusion. A
service with a good mean and a fat left tail needs the second, and an
intervention aimed at the first will not help it.
\end{reading}

\begin{workedexample}[apb-driftdiff]{Drift and diffusion, separated numerically}
Start a clinic at a Gaussian density on $x\in[0,10]$ with $401$ cells, centred
at $5.0$ with variance $0.25$, and take the protective level to be $4.5$. At
the start, $15.3$ per cent of the clinic sits below it. Run the
Fokker--Planck step of Proposition~\ref{prop:B_mass} for $50$ backward-Euler
steps of $dt=0.02$, that is to $t=1$, under three policies.

Doing nothing, with only the baseline diffusion $D=0.05$, leaves the mean at
$5.00$, raises the variance from $0.25$ to $0.30$ and pushes the tail to
$17.5$ per cent. A drift policy, $\mu=0.5$ with the same $D$, moves the mean to
$5.50$ and leaves the variance at $0.32$; the tail collapses to $3.6$ per cent.
A variability policy that instead raises $D$ to $0.8$ with no drift leaves the
mean at $5.00$ and takes the variance to $1.05$; the tail rises to $30.8$ per
cent.

\begin{wecode}[Python]
import numpy as np
from hpv.dynamics import fokker_planck_step

x = np.linspace(0.0, 10.0, 401); dx = x[1] - x[0]
base = np.exp(-((x - 5.0) ** 2) / 0.5); base /= base.sum() * dx
mean = lambda r: float((x * r).sum() * dx)
var = lambda r: float(((x - mean(r)) ** 2 * r).sum() * dx)
tail = lambda r: float(r[x < 4.5].sum() * dx)

def run(mu, D, steps=50, dt=0.02):
    r = base.copy()
    for _ in range(steps):
        r, _ = fokker_planck_step(r, x, mu, D, dt)
    return r

zero = np.zeros_like(x)
runs = {"start": base,
        "no policy": run(zero, np.full_like(x, 0.05)),
        "drift 0.5": run(np.full_like(x, 0.5), np.full_like(x, 0.05)),
        "diffusion 0.8": run(zero, np.full_like(x, 0.8))}
for name, r in runs.items():
    print(f"{name:<14} mean {mean(r):.3f}  var {var(r):.3f}  tail<4.5 {tail(r):.3f}")
\end{wecode}

\weanswer{Drift moves the mean $5.00\to5.50$ and leaves the variance at $0.32$;
diffusion holds the mean at $5.00$ and moves the variance $0.30\to1.05$. The
tail below the protective level goes to $0.036$ under drift and $0.308$ under
diffusion, so the two policies are opposite in effect on the quantity that
matters. Asserted by
\texttt{test\_worked\_example\_drift\_and\_diffusion\_in\_numbers} and
\texttt{test\_fokker\_planck\_separates\_drift\_from\_diffusion}.}
\end{workedexample}
\checkedby{tests/test_dynamics.py::test_proposition_fokker_planck_conserves_mass}

\begin{remark}[Boundary conditions and interpretation]
Reflecting boundaries model a closed cohort in which patients neither enter nor
leave the biomarker range; absorbing boundaries would model an exit event (for
example death at $x=0$). The stationary density, where it exists, is the
long-run population distribution of disease states, and its variance and tail
mass are the quantities a policy would seek to shift.
\end{remark}

For a cohort of size $N$ the coupled system tracks individual states $X_i(t)$
with summaries $\overline X(t)=\tfrac1N\sum_i X_i(t)$ and
$\operatorname{Var}[X](t)=\tfrac1N\sum_i(X_i-\overline X)^2$; heterogeneity enters
through distributions of the kinetic parameters.

The three layers describe how a disease marker moves. They say nothing yet
about the conceptual score of Chapter~\ref{ch:theory}, which was defined at a
single moment. Connecting them means letting that score move too.

\section{A Time-Varying Form of the Conceptual Model}
\label{sec:B_timevarying}

\begin{definition}[Time-varying biomedical term]
\label{def:B_timevarying}
Let $L(t)=\phi(X(t),\overline X(t),\text{comorbidities},\text{age})$ with $\phi$
differentiable and increasing in $X$. The time-resolved metric is
\begin{equation}
H_{\text{potential}}(t)=L(t)\,Q(t)\,S(t)-B(t).
\label{eq:B_timevarying}
\end{equation}
\end{definition}

\begin{proposition}[Differentiability]
\label{prop:B_differentiable}
If $L,Q,S,B$ are differentiable in $t$, then $H_{\text{potential}}(t)$ is
differentiable, with
$\dot H = \dot L QS + L\dot Q S + LQ\dot S - \dot B$.
\end{proposition}
\begin{proof}
Products and differences of differentiable functions are differentiable; the
product rule applied to $LQS$ gives the stated derivative.
\end{proof}

\begin{reading}[Reading $\dot H$ as a rate of change of capability]
The score differentiated here is the heuristic $(LQ)S-B$, which
Chapter~\ref{ch:theory} demotes and Chapter~\ref{app:crf} replaces. It is used
in this appendix because it is the simplest object that can be made to move in
time, and the structural point below carries over unchanged to the
realisation function.

The derivative $\dot H=\dot LQS+L\dot QS+LQ\dot S-\dot B$ has one term per
thing that can change, and the coefficient on each says how much that change is
worth in this patient's current circumstances.

The middle terms carry the argument. The value of improving quality, $\dot Q$,
is multiplied by $LS$: it is worth more in someone with years ahead of them and
circumstances that let them use those years. The value of improving social
conditions, $\dot S$, is multiplied by $LQ$: worth more in someone whose
biology still has something to give.

For Profile~B this is the difference between two funding decisions taken at
different ages. At age $58$, with $L$ still substantial, a change in $S$ or $B$
buys a lot. The same change at age $85$ buys much less because $L$ has shrunk.
The model states that trade openly; it claims no more than that.
\end{reading}
\checkedby{tests/test_dynamics.py::test_time_varying_score_is_differentiable}

This is a theoretical generalisation of the static model of
Chapter~\ref{ch:theory}; it is not estimated or validated here. Simulating
$X(t)$ over time would yield time-resolved trajectories for individuals and
cohorts.

What the appendix has established, and what it has only proposed, are worth
separating before the disease-specific version is attempted.

\section{Summary}
\label{sec:B_summary}
This appendix set out a disease-agnostic ODE--SDE--PDE architecture: a
deterministic backbone with a unique stable equilibrium
(Propositions~\ref{prop:B_existence}--\ref{prop:B_equilibrium}), a stochastic
extension that reduces to the backbone as noise vanishes and whose mean obeys the
deterministic law (Propositions~\ref{prop:B_reduction}--\ref{prop:B_consistency}),
a conservative population-level Fokker--Planck layer
(Proposition~\ref{prop:B_mass}), and a differentiable time-varying form of the
conceptual score (Proposition~\ref{prop:B_differentiable}). The next appendix
specialises this framework to haemophilia; Appendix~\ref{app:computation} gives
its numerical implementation.

Every proposition above is checked automatically. Nine of the fourteen tests in
\texttt{verification/tests/test\_dynamics.py} assert the Lipschitz constant, the
convergence of nearby trajectories, agreement between the closed-form
equilibrium and the integrator, the vanishing-noise limit, the affine-drift
mean, mass conservation under reflecting boundaries, the separation of drift
from diffusion, and the product rule for the time-varying score.
Appendix~\ref{app:verification} prints the result of the run that produced this
document, and records the sign error in the Fokker--Planck assembly that the
drift test caught while the conservation test passed.
